% RQ4 Background: Treasury Management

\subsection{DAO Treasuries in Practice}

\subsubsection{Treasury Composition}

Typical DAO treasuries contain:
\begin{itemize}
    \item \textbf{Native tokens}: Governance tokens (often majority of treasury by value)
    \item \textbf{Stablecoins}: USDC, DAI for operational stability
    \item \textbf{Protocol revenue}: Accumulated fees, yields
    \item \textbf{External assets}: Investments, diversification holdings
\end{itemize}

Native token concentration creates correlation risk: treasury value drops precisely when the protocol struggles.

\subsubsection{Treasury Uses}

DAOs spend treasury on:
\begin{itemize}
    \item Development grants and bounties
    \item Operational expenses (infrastructure, services)
    \item Marketing and growth initiatives
    \item Liquidity provision and incentives
    \item Strategic investments and acquisitions
\end{itemize}

\subsubsection{Historical Challenges}

Notable treasury crises include:
\begin{itemize}
    \item Bear market depletion (2022): Multiple DAOs saw 80-90\% treasury value loss
    \item Unsustainable spending: DAOs approving grants that exceeded sustainable burn rate
    \item Concentration losses: Treasuries holding single assets that collapsed
\end{itemize}

\subsection{Existing Treasury Practices}

\subsubsection{Compound}

Compound maintains treasury primarily in COMP tokens with limited diversification. No formal buffer policy; spending governed by individual proposal votes.

\subsubsection{Uniswap}

Uniswap's treasury exceeds \$1B but is almost entirely UNI tokens. The Uniswap Grants Program manages a separate allocation for grants.

\subsubsection{Optimism}

Optimism employs structured treasury management with dedicated allocations for different purposes (ecosystem fund, governance fund, etc.) and vesting schedules.

\subsubsection{MakerDAO}

MakerDAO represents relatively sophisticated treasury management with stability reserves, surplus buffer, and systematic approach to capital deployment.

\subsection{Theoretical Foundations}

\subsubsection{Corporate Finance Parallels}

Treasury management in DAOs parallels corporate cash management:
\begin{itemize}
    \item Cash buffer theory: Firms hold precautionary balances for uncertainty
    \item Target cash models: Optimal cash depends on transaction costs, opportunity cost, volatility
    \item Liquidity management: Tradeoff between idle cash and operational needs
\end{itemize}

\subsubsection{Reserve Adequacy}

Central banks and insurers use reserve adequacy frameworks:
\begin{equation}
\text{Reserve ratio} = \frac{\text{Liquid reserves}}{\text{Expected outflows}}
\end{equation}

DAOs can adapt these frameworks for treasury policy.

\subsection{Related Work}

Limited academic work addresses DAO treasury management specifically. \citet{faqir2021cadcad} modeled tokenomic dynamics that affect treasury value. DeepDAO \citep{deepdao2023analytics} provides treasury analytics across DAOs. Our simulation framework enables systematic study of treasury policies.
